\section{Part 1：Scaling laws 的历史和数据 scaling}

前面的工程问题需要一个可信前提：小规模实验是否真的能揭示大规模趋势。本部分先从统计学习中的 sample complexity 和早期 NLP 经验出发，建立“误差随数据量按幂律下降”的证据链，并区分理论上界、经验拟合与可用于预算决策的预测模型。


\singleslide{slides-images/slide-005.jpg}{Part 1 路线：把 data scaling 看作经验 sample complexity，并追溯神经网络早期 scaling 研究。}{5}

这页规定历史部分的用途：理论 sample complexity 给出“需要多少数据”的上界或速率，经验 scaling law 则测量真实模型在真实分布上的 loss。二者都讨论随数据增长的误差衰减，却回答不同问题；后文会用差异解释为何 power law 既熟悉又令人惊讶。

\begin{knowledgebox}{讲义补充：源 Slide 5 的 Part 1 不是只讲历史}
这一部分先把 scaling laws 放进学习理论和经验 NLP 的背景：理论上早就有 sample complexity 和 rate 的概念，但它们往往是上界；现代 scaling laws 更关注实际模型、实际数据和实际 loss 曲线是否可预测。
\end{knowledgebox}


\singleslide{slides-images/slide-006.jpg}{经典 sample complexity 与 rates：有限假设集和光滑密度估计给出上界，但不是实际 realized loss。}{6}

Slide 6 对比理论与工程测量。VC/finite-hypothesis 分析常保证误差不超过某个量，非参数统计给出随 $n$ 的收敛速率；这些结果不直接告诉某个 Transformer 在某数据集上的具体 cross-entropy。Scaling experiments 的价值正是补上常数、指数与真实 regime。

\begin{knowledgebox}{讲义补充：源 Slide 6 的理论上界与实际 loss}
理论 sample complexity 常告诉我们达到某个 error 需要多少样本，或者 error 上界如何随 \(n\) 衰减。但 frontier LMs 关心的是实际验证 loss 在真实训练分布上的可实现值。Scaling laws 更像工程经验曲线，不等同于严格泛化上界。
\end{knowledgebox}


\singleslide{slides-images/slide-007.jpg}{1993 年早期 data scaling law：在现代深度学习之前就观察到随数据规模的规律性改进。}{7}

早期论文说明 log-log 线性不是 LLM 时代才出现的营销图形。它把模型性能与训练样本数联系成可外推函数，开启了“用较小实验估计更多数据价值”的经验路线。现代神经 scaling 的新意在于规律跨任务、模型与更大数量级持续成立。

\begin{knowledgebox}{讲义补充：源 Slide 7 的历史意义}
这页说明 scaling 的想法并非 LLM 时代突然出现。早期研究已经观察到更多数据与性能之间的规律关系。LLM 时代的新东西是规模更大、曲线更平滑、投资更昂贵，因此这些规律变成了核心工程工具。
\end{knowledgebox}

\begin{figure}[H]
\centering
\includegraphics[width=0.86\textwidth]{slide-008.jpg}
\caption{Slide 8：Banko and Brill 2001 观察到随数据增长的 log-linear scaling。}
\figsource{CS336 Spring 2026 Lecture 9 官方幻灯片。}
\end{figure}

\begin{knowledgebox}{读图：Slide 8 的横纵轴直觉}
这类图通常横轴是数据量，纵轴是错误率或性能指标。重点不是某个点，而是随着数据增加，性能按简单函数改善。早期 NLP 的经验已经提示：更多数据在相当长范围内有可预测收益。
\end{knowledgebox}


\singleslide{slides-images/slide-009.jpg}{早期 data scaling 的函数形式检验：下游性能与数据量呈 power-law 关系。}{9}

Slide 9 的重点不是某个任务的精确指数，而是模型选择：指数、对数或饱和曲线在有限区间都可能拟合，只有跨数量级实验与 held-out extrapolation 才能区分。Power law 在 log-log 坐标上近似直线，使指数成为可比较的“数据效率”摘要。

\begin{knowledgebox}{讲义补充：源 Slide 9 应该比较函数形式}
这页的教学点是：scaling law 不是先验假设某条线，而是比较不同函数形式对实验点的拟合和外推能力。Power law 之所以重要，是因为在 log-log 坐标中会变成近似直线，便于拟合和解释。
\end{knowledgebox}

\begin{figure}[H]
\centering
\includegraphics[width=0.86\textwidth]{slide-010.jpg}
\caption{Slide 10：Hestness et al. 2017 在机器翻译、语言模型、语音等任务上看到可预测 scaling。}
\figsource{CS336 Spring 2026 Lecture 9 官方幻灯片。}
\end{figure}

\begin{knowledgebox}{读图：Slide 10 的跨任务稳定性}
Hestness 的重要性在于跨任务：MT、LM、Speech 都呈现规律 scaling。若 scaling 只在单个 benchmark 上出现，它只能算偶然经验；跨任务出现说明某些统计或优化机制具有普遍性。
\end{knowledgebox}


\singleslide{slides-images/slide-011.jpg}{Hestness 等工作的前瞻结论：emergence、compute scaling，以及速度与准确率的统一趋势。}{11}

这页把 data scaling 扩展到 compute 与 qualitative behavior。所谓 emergence 可能来自指标阈值或连续 loss 改进在离散任务上的投影；“speed = accuracy”则提示更多 compute 不只提高最终质量，也可能缩短达到目标精度所需时间。后文因此同时拟合 data、model 与 compute。

\begin{knowledgebox}{讲义补充：源 Slide 11 连接到今天的 emergence 讨论}
这页提示，许多今天看似新的问题，如能力涌现、计算速度和准确率 trade-off，早期 scaling 工作已经触及。Scaling laws 不只预测 loss，也帮助我们判断某类能力是否可能通过 brute force scale 获得。
\end{knowledgebox}

\subsection{本章小结}

Part 1 的历史线索说明：scaling laws 不是 LLM 独有魔法，而是数据、模型和计算规模增加时经常出现的经验规律。LLM 的特殊性在于这些规律足够平滑、成本足够高，因而成为训练决策的基础设施。

